\documentclass{ximera}

\begin{document}

    \title{Lecture 21}
    
    \begin{abstract}  
    Outline: \\
    - Interior Extremum Theorem \\
    - Darboux's Theorem 
\end{abstract}

\maketitle


\begin{theorem}{(Interior Extremum Theorem)}
    Let $f$ be differentiable on $(a,b)$.  If $f$ attains a max/min at some point $c \in (a,b)$, then $f'(c) = 0$.
\end{theorem}
\begin{proof}
    Since $c \in (a,b)$, we can construct $x_n \to c$ with $x_n < c$, and $y_n \to c$ with $y_n > c$ for all $n \in \mathbb{N}$.  Notably,
    \begin{align*}
        f'(c) = \lim\limits_{n\to \infty} \frac{f(y_n) - f(c)}{y_n - c}, 
    \end{align*}
    but each difference quotient is nonpositive since $f(c)$ is a maximum!  Similarly, 
    \begin{align*}
        f'(c) = \lim\limits_{n\to \infty} \frac{f(x_n) - f(c)}{x_n - c}, 
    \end{align*}
    but each difference quotient is nonnegative since $f(c)$ is a maximum!  Thus, by the Order Limit Theorem, $0 \leq f'(c) \leq 0$, or $f'(c)=0$.
\end{proof}

\begin{exercise}
    Prove Darboux's Theorem:
    \begin{theorem}
        If $f$ is differentiable on an interval $[a,b]$, and if $\alpha$ satisfies $f'(a) < \alpha < f'(b)$ (or $f'(a) > \alpha >f'(b)$), then there exists a point $c \in (a,b)$ where $f'(c)= \alpha$.
    \end{theorem}
    \textit{Hint:} Define a new function $g(x) = f(x) -\alpha x$, so that $g'(x) = f'(x) - \alpha$.  Then, $g'(a) < 0$ and $g'(b) > 0$.  Show that there exists a point $x \in (a,b)$ such that $g(a) > g(x)$, and a point $y \in (a,b)$ where $g(y) < g(b)$.
\end{exercise}
\vspace{-0.3in}
\textcolor{red}{
\begin{proof}
    Define a new function $g(x) = f(x) -\alpha x$, so that $g'(x) = f'(x) - \alpha$.  Specifically, since $g'(a)< 0$, then there exists a $\delta > 0$ such that the difference quotient
\begin{align*}
    \frac{g(x) -g(a)}{x-a} < 0 
\end{align*}
for some $0 < x-a < \delta$.  Consequently, since $x-a$ is positive, $g(x) - g(a) < 0$, or $g(x) < g(a)$.  An analogous proof holds to show that there exists $g(y) < g(b)$.  Moreover, since $g$ is continuous, $g([a,b])$ is a compact set, and consequently contains its maximum and minimum.  The minimum cannot be at the endpoints based on our argument that $g(y) < g(b)$ and $g(x) < g(a)$, and therefore it must be in the interior at some point $c \in (a,b)$.  Then, we can apply the Interior Extremum Theorem and state that there exists a point $c\in (a,b)$ where $g'(c) = 0,$ which implies $f'(c) = \alpha$.
\end{proof}
}

Note that this is kind of like the Intermediate Value Theorem.  However, it is not an immediate application of this, as there exist functions which are differentiable, but their derivative is not continuous.  Our example $f(x) = x^2\sin(1/x)$ has derivative $f'(x) = 2x\sin(1/x) - \cos(1/x)$, and $\cos(1/x)$ is not continuous at $x=0$.


\end{document}
